\documentclass{article}
\usepackage[utf8]{inputenc}
\usepackage[T1]{fontenc}
\usepackage[a4paper,total={6in,8in}]{geometry}
\usepackage{amsmath,amssymb,amsfonts}
\usepackage{graphicx}
\usepackage{hyperref}
\usepackage{enumitem}
\usepackage{listings}
\usepackage{xcolor}
\DeclareMathOperator{\Var}{Var}
\DeclareMathOperator{\Cov}{Cov}
\providecommand{\tightlist}{\setlength{\itemsep}{0pt}\setlength{\parskip}{0pt}}
\lstset{basicstyle=\ttfamily\scriptsize,breaklines=true,columns=fullflexible,
  keepspaces=true,showstringspaces=false,frame=single}
\hypersetup{colorlinks=true,urlcolor=blue,linkcolor=blue}
\setlength{\emergencystretch}{2em}

\title{Math 126 Problem 18: Finding Cliques}
\author{} % Put your name here.
\date{Spring 2027}
\begin{document}
\maketitle

\hypertarget{submission-instructions}{%
\subsection*{Submission Instructions}\label{submission-instructions}}

Please upload your work to Gradescope by \textbf{11:59pm on Monday, April 12, 2027 (Pacific time).}

\begin{itemize}
\item
  You are encouraged to discuss ideas and work with your classmates. However, you \emph{must acknowledge} your collaborators at the top of each solution on which you collaborated with others and you \emph{must write} your solutions independently.
\item
  Your solutions to theory questions must be written legibly, or typeset in LaTeX or markdown.
\item
  I recommend that you write your solutions in a Python notebook using Google Colab.
\item
  You should submit your solutions as a \emph{single PDF} via the assignment on Gradescope.
\end{itemize}

\textbf{Grading:} The point of the problem set is for \emph{you} to learn. To this end, I hope to disincentivize the use of LLMs by \emph{not} grading your work for correctness. Instead, you will grade your own work by comparing it to my solutions. This self-grade is due the Friday \emph{after} the problem set is due, also on Gradescope.

\noindent Original notebook: \href{https://colab.research.google.com/drive/19daxPN0hhRwvyQ3BLb0z7GdGOTqvOVRK?usp=sharing}{Google Colab}.

% Source notebook cell 2
\hypertarget{finding-cliques}{%
\subsection*{Finding Cliques}\label{finding-cliques}}

A common task in data mining is to identify large \emph{cliques} in a graph. For example, in social networks, large cliques can be indicators of fraudulent accounts or networks of accounts designed to promote certain content. In this problem, we consider a spectral heuristic for finding a large clique based on the top eigenvector of the graph adjacency matrix \(\mathbf{A}\):

\begin{itemize}
\tightlist
\item
  Compute the leading eigenvector \(\mathbf{v}_1\) of \(\mathbf{A}\).
\item
  Let \(i_1, \ldots, i_k \in \{1, \ldots, n\}\) be the indices of the \(k\) entries in \(\mathbf{v}_1\) with the largest absolute value.
\item
  Check if nodes \(i_1, \ldots, i_k\) form a \(k\)-clique.
\end{itemize}

We will analyze this heuristic on a natural random graph model. Specifically, let \(G\) be an Erdős-Rényi random graph: we start with \(n\) nodes, and for every pair of nodes \((i,j)\), we add an edge between the pair with probability \(p < 1\). To simplify the math, also assume that we add a self-loop at every vertex \(i\) with probability \(p\). Then, choose a fixed subset \(S\) of \(k\) nodes to form a clique. Connect all nodes in \(S\) with edges and add self-loops. We will argue that, for sufficiently large \(k\), we can expect the heuristic above to identify the nodes in the clique.

\hypertarget{part-a-expected-adjacency}{%
\subsection*{Part A: Expected Adjacency}\label{part-a-expected-adjacency}}

Let \(\mathbf{A}\) be the adjacency matrix of a random graph generated as above. What is \(\mathbb{E}[\mathbf{A}]\)? Then write \(\mathbb{E}[\mathbf{A}]\) in terms of its eigenvalues and eigenvectors (you don't need to find the exact algebraic expressions for these!).

\textbf{Hint:} Argue that, up to multiplying by a constant, any eigenvector \(\mathbf{v}\) corresponding to a non-zero eigenvalue must have \(v_i = 1\) for all \(i \notin S\) and \(v_i = \alpha\) for all \(i\in S\), where \(\alpha\) is a constant.

\hypertarget{part-b-random-adjacency}{%
\subsection*{Part B: Random Adjacency}\label{part-b-random-adjacency}}

To prove the algorithm works, it is possible to use a matrix concentration inequality to argue that the top eigenvector of \(\mathbf{A}\) is close to that of \(\mathbb{E}[\mathbf{A}]\). Instead of doing that, let's verify things experimentally. Generate a graph \(G\) according to the prescribed model with \(n = 900\), \(k = |S| = 30\), and \(p = 0.1\). Compute the top eigenvector of \(\mathbf{A}\) and look at its 30 largest entries in magnitude. What fraction of nodes in the clique \(S\) are among these 30 entries? Repeat the experiment 100 times and report the average fraction recovered.

\end{document}
